% Source: https://artofproblemsolving.com/community/c3h2363616_tesseractexe
% Stable identity: aops-topic-2363616-q01
A tesseract (a four-dimensional cube) can be described as the set of points $(x_1,x_2,x_3,x_4)$ with $0\le x_i\le1$ for each $i$. Its vertices are the points whose coordinates are all $0$ or $1$; its edges, square faces and cubical cells are the pieces obtained by letting $1$, $2$ or $3$ of the coordinates vary between $0$ and $1$ while each of the others is fixed at $0$ or at $1$. How many
\begin{enumerate}
\item[(a)] cells,
\item[(b)] square faces,
\item[(c)] edges,
\item[(d)] vertices
\end{enumerate}
does a tesseract have?
